\section{第一章　坠入兔窟}

“模型的准确率上不去一定不是我的问题。”玲嘟囔着，断开手机上的远程连接，下了出租车。如果说生活是由一连串败坏心情的日子组成，那么今天的令人沮丧的程度显然是个离群值（此处略去统计学检验）；而玲宁愿不要这样的每一天的连续性。训练了一晚上的AI模型的各项量标相比昨天深夜（或者是今天凌晨？）她离开时，竟然没有提升，反而下降了若干个百分点。当她被不同于平日的闹钟叫醒（又一个离群值），睡眼朦胧地坐上出租车，在说服自己已经回不去被窝的同时打开手机连接训练集群，看到不争气的准确率在图表中上下振荡，又一想到自己被迫早起来这儿只是为了参加一个狗屁研讨会，玲的心跳似乎都变得无力了。

“好吧，至少让我欣赏欣赏同行们的‘杰作’。”玲拖着步子走向电梯。这里是C座1层；那么她要去的C座4层应该就在乘坐电梯向上的第3层（这是不言自明的；对此有疑虑的，建议复习幺半群上的运算或者建筑的图纸）。当电梯磨磨蹭蹭地降下来并打开门时，玲已经无聊到开始思考电梯带来的加速度事实上与空间扭曲等效。她没好气地走进电梯，拼命连按关门键。要是有人现在想要走进来，电梯门保准会夹断他的腿，玲心想。

然而伸进来的不是腿，而是一只耳朵。“等等！”玲听到这声大喊，反应过来时，发现自己的手已经在按开门键了。不等玲开始思考条件反射和无条件反射的区别，电梯外就走进一位穿着立领衬衫戴着白手套的绅士。“你差点夹到我的耳朵！”玲没有听他的抱怨，而是打量着他的手套，想他为什么做完实验不摘手套（这违反生物安全规定！），而后发现他戴的是一双皮手套，又开始想这位绅士来这儿干什么。难道是玲的打扮太过随意了？

那位绅士没有按电梯按钮（看来他也是去同一楼层；或者是他忘记按了），只是盯着手机上的时钟嘀咕着：“哎哟！哎哟！来不及了！我要迟到了！”当然，戴着手套操作手机是很滑稽的；玲不知道他是如何能够隔着手套让触摸屏响应他的飞快点按的。

“早着呢，研讨会还没开始。”刚说出口，玲就后悔了。首先，她并不知道这位绅士是不是和他一样来参加研讨会的；即使要去的楼层一样，也不能说明什么。其次，对方着急不着急，和自己又有什么关系呢？

“那是因为你不用在会上发言。”绅士白了她一眼。这是什么意思？他是怎么知道玲不用发言的？难道在这某不知名研究所举办的某不知名会议上发言还有着装要求？另外，翻白眼是否符合贵族礼节？是否不应该再称呼他为“绅士”了？

好在4层终于到了。看着“绅士”一蹦一跳地跑出电梯，玲感到有点滑稽。玲跟着走出了电梯。她面前是一条很长的走廊，走廊两侧分布着大小不一的门（现代建筑艺术吧？玲不懂这个）。玲现在懒得思考了，跟着“绅士”走进一扇门，里面正是会场。

会场里闹哄哄的；这也正常。与会者们聚成一个个小团体，热烈地交流着，也不知道他们聊的是学术话题还是学术八卦（也许后者也是前者的一个子领域）。玲找了一个角落里的空位置坐下，托着头等待会议开始。诶，哪个缺德的把会议时间定在早晨；以后应该开个关于“回笼觉有益身心健康”的研讨会（尤其是对于像她这样半夜还在训练模型的猝死预备人群来说）。

“大伙儿安静一下，”主持人上台了，“会议马上就要开始了……”

（此处省略毫无营养的主持人长篇开场辞，包括但不限于：介绍和感谢到场嘉宾并要求鼓掌，不知道从网上哪儿抄来的AI发展历程，纯粹臆想的前沿和挑战，等等等等。）

第一位发言者踩着高脚凳出现在台上。玲坐得远，实在看不清他的相貌（也许不全是坐得远的问题）。他开始作一个“使用RWKV\footnote{Receptance Weighted Key Value，线性加权门控循环网络（自拟译名）}代替Transformer\footnote{自注意力变换器（自拟译名）}的可行性报告”。当他开始推导公式时，玲都快睡着了。当她再次清醒过来时，已经到了问答环节（或者说激烈的争论环节）。

“这有什么意义，还不如再加大点键值缓存\footnote{KV cache。“键值缓存”为自拟译名}。”

“既然你说它这么好，那主流大模型怎么没一个用它呢？”

如此种种。玲实在提不起兴趣；她不是研究理论方向的。最后还是靠主持人强行打断了，“时间差不多了！”

第二位发言者是个秃头，眼睛都快长到头顶上去了，还戴着一副不成比例的黑框眼镜。他上来就说所有AI研究者都应该拜康德为祖师爷，接着就开始大谈康德哲学和AI的关系：

“康德对Vorstellung\footnote{表象}和Ding an sich\footnote{物自体}的区分是基础性的。我们永远无法直接把握现实的原貌，就像AI模型仅仅能认识到数据集的采样，而不能认识到基准真相\footnote{Ground truth}本身；就连基准真相也不过是对现实世界的采样……

“AI模型的认知受Formen der Anschauung\footnote{直观形式}制约：比如，卷积神经网络\footnote{Convolutional Neural Networks (CNN)}明确地嵌入了一种空间先验，而循环神经网络\footnote{Recurrent Neural Networks (RNN)}具有时间形式的内在结构……

“神经网络的结构、层类型、激活函数和超参数正如康德所说的Kategorien\footnote{范畴}：它们决定了AI模型以什么样的方式去理解它的输入……

“注意力机制动态地确定输入序列的不同部分如何相互关联以产生输出；这种机制充当一种复杂的Schematismus\footnote{图式}……

“模型训练中的最优化过程，比如梯度下降，正是推动apriorische Synthesis\footnote{先验综合}的动力……

玲听不下去了。且不说搬出一个两百多年前的老古董，用拙劣的类比生搬硬套有什么实际意义，单说这位发言者一会儿汉语一会儿德语就够令人难受的了。更好笑的是，他还是用英语的发音念德语单词。“谁邀请人文社科的臭蛤蟆来的！这像话吗？”玲暗暗不爽。

“臭蛤蟆”终于讲完了，好在没有拖到超时。玲环视四周（当然，她是坐在角落的，所以严谨来说应当是四周的四分之一）；看起来和她类似想法的不在少数。大家窃窃私语，而没人愿意提问。“臭蛤蟆”尴尬地在台上站了会儿；好在最后主持人给他解了围。

接着上台的竟是先前玲在电梯里遇到的“绅士”。他彬彬有礼地鞠了一躬（这才像话嘛！难道绅士都是人前一套人后一套？），耳朵碰了碰讲台，开始他的发言：

“解释并调优AI模型一直以来都是一项困难的工作；其阻碍的根本原因，在于AI模型对于我们研究者来说是一个黑箱：我们仅仅能够从外部，通过各种量标去观测它；而对其内在机理却一无所知。这导致了一个结果：我们常常在AI模型的训练中观察到匪夷所思的突变，比如说，涌现、崩溃，却不知道其所以然……

玲不由地点了点头。总算来了个她感兴趣的；这次会议也算没白来嘛！

“我们最近的工作的贡献就在于，我们创新性地发明了一种全新的方法，用以从内部探究AI模型的机理，那就是——使用冷冻电镜对AI模型进行成像。我们使用冷冻电镜技术通过超低温固定AI模型的动态构型，解析其内部语义空间、知识存储路径及逻辑决策节点的微观结构。通过这样的技术革新，我们得以在张量级别研究AI模型的结构和动力学，进而为AI的可解释性提供了坚实的基础……

等会儿？玲甚至怀疑自己听错了（要是她有这位“绅士”一样长的耳朵就好了）。什么玩意儿？冷冻电镜应该说的是冷冻电子显微术吧？这和AI有半毛钱关系吗？不对不对，首先，AI模型怎么会有物理实体？

台上的“绅士”继续煞有介事地讲着，甚至展示了好几张看起来像那么一回事的影像和图表。玲期待会有脑子正常的听众起哄；然而没有；大家看起来都听得很认真。

在玲持续怀疑是自己的脑子坏掉了还是其他人的脑子都坏掉了的时候，发言已接近尾声。其他人肯定是憋着质疑等问答环节再大加攻伐，玲这样想。

问答环节确实有质疑，但完全不是玲所想的那样：

“请问你们为什么要选用冷冻电子显微术对AI模型进行单颗粒分析呢？我觉得使用冷冻电子断层扫描术更能捕捉AI模型的整体结构吧？”

“请问你们的技术的分辨率真能达到张量级别吗？我不太相信。你们是不是做了什么预处理或者后处理？”

“请问你们的成像分析软件开源了吗？”

……

好吧，一定是玲疯了，或者是这个世界疯了。玲一刻也不想待下去了。好在接下来就是中场休息的茶歇时间（或者说，蹭吃蹭喝时间）。玲想起自己来的时候还没吃过早饭，就想赶紧吃点然后开溜。

玲在茶歇上的煎饼摊点了一份馃子，啃了起来。（等会儿，茶歇上怎么会有煎饼馃子？算了，今天的怪事已经够多了；若是一件一件地尝试去思考，怕是再用几个大语言模型也想不过来。）正吃着，玲看到几个人围着一个人聚了起来。吃东西的时候也没法干其他有意思的事情，于是玲便无聊地上前看看。

人群中间是一位戴着眼镜的人，一副书呆子扮相。（要是在平常，这就是研究人员常见的打扮；不过今天也不能算作平常。）看起来他正在吹嘘他新写的短篇小说（所以说，人文社科的人来得太多了！），名字叫什么《延迟中的幽灵》。他说这是一篇探讨AI的充满隐喻的小说。

玲最讨厌这种末流写手用什么奇奇怪怪的被称之为“文学性诠释”的比喻来扭曲他们研究者的成果了。她无不讥讽地问道：“那么请问阁下，AI能理解你的关于AI的充满隐喻的‘杰作’吗？”

四眼狗立刻紧张起来；玲甚至能看到他的镜片都在流汗。他支支吾吾地自言自语：“这立刻产生了一个自我指涉的悖论……”

看到他答不上来，玲心中升起一股暗爽；虽然这种通过贬低他人产生的感受立刻使她背上一种违背道德的负罪感（什么道德？康德的伦理学吗？玲之前有为自己立过法吗？玲回忆起会上“臭蛤蟆”的发言）。她忽然发觉手上的煎饼已经吃完了，并且感到自己很渴（这是显而易见的；煎饼馃子普遍地来说是又干又咸的）。她便留下人群，自顾自地去找饮料喝。

玲找到一瓶牛奶，便立刻打开盖子往干涸的喉咙里灌。等她喝完了，才发觉味道不太对。她看了一眼，发现玻璃瓶的标签上写的是“兔子奶”。“那怪不得，”玲心想，“我先前对于它是牛奶的判断是完全的独断论。”玲又想了想，兔子是哺乳动物吗？确实是的。但拿兔子产奶真的符合经济利益吗？也许是养殖技术有新的进步吧。

玲感到自己越来越困了。玲想，其一是因为自己今天确实睡眠不足；其二是因为刚吃饱导致血液富集向胃肠。正当玲想要赶紧开溜回家补觉时（昨天的模型都训练成这副鸟样了，今天也不知道该怎么调整；事已至此，先睡觉吧），“绅士”一边看着手机，一边急匆匆地从她身边跑过去（就像他来时跑出电梯那样急），一边还忧心忡忡地嘀咕着：“不好了！不好了！昨天的模型做完冷冻电镜之后怎么梯度消失了！这下肯定要被大老板骂了！”

玲还是想知道是自己疯了还是其他人疯了，于是尝试和“绅士”搭话，想要问问他的“冷冻电镜AI模型成像”究竟是不是愚人节笑话。她在“绅士”身后追着，拐过走廊里三四个拐角；可最终他还是跑没影了。玲这下知道，“守株待兔”或许是有一定道理的。

还是赶紧回去吧。说不定睡一觉之后这个世界就会恢复正常了。玲这样想着，走进前方最近的一部电梯。她按下1层的按钮，电梯门缓缓关上。

玲靠在电梯的墙壁上，努力让自己保持清醒。至少到家之后再睡吧！不过现在怎么这么困；这种困意是玲从来没有过的。不知过了多久，当玲重新撑开眼睛时，她发现电梯还在运行，自己还没到1层。玲看了看显示屏上的数字：现在是2.50层；这个数字而后立刻接着跳动起来。2.49、2.48、2.47……

好嘛！原来不是自然数的幺半群，而是有理数的幺半群；不过它在更严格意义上是否也满足群的条件，这还有待考察。玲思考着这意味着什么。看来自己确实处于奇异事件中。如果按照某基金会的叙事风格，现在自己应该被列为已失踪的D级人员了。

“如果我训练的模型的损失值也能像这样持续下降就好了。”玲对自己开玩笑，尝试缓解自己的紧张和恐惧。在漫长的等待之后（这似乎比盯着模型，看它一步一步训练还漫长），电梯终于停下，1层终于到了。电梯门缓缓打开。

“现在即使门外是异世界也不奇怪了。”玲走出电梯，发现自己正站在天台上。

玲对自己向下坐电梯却来到天台这一情况提出两种假说：

\begin{enumerate}
\item 电梯向下穿过了整个地球，带她来到了地球的另一端。且不论这在工程上是否能实现，这一假说立刻产生了一个无法回避的问题：为什么在穿越地球的同时，玲没有察觉到重力的减轻和逆转？
\item 有大质量的天体接近地球，扰乱了地球的引力场。根据广义相对论的等效原理，这确实有可能在电梯启动时让玲感到失重而在电梯停止时让玲感到超重，从而使玲产生了电梯在下行的印象；而电梯实际上则是在上行。
\end{enumerate}

玲认为第二种假说更加符合逻辑。不过现在也没机会多想了，因为有更紧迫的情况。玲看到一个身穿纯白衣服的女孩站在天台的栏杆外，向前举起双臂。

“危险！”玲赶紧跑过去，大喊着。现在应不应该大喊呢？或许不应该吧；女孩受到惊吓可能会导致更糟的情况发生。而且也不应该喊“危险”而应该喊“等一下”，因为如果女孩是要寻短见，危不危险便不在她的考量范围之内；毋宁说，越危险越能够达成目的。不过玲都已经喊出口了。

女孩的双手放开了。玲已经十分接近女孩了；她这时才看清之前在女孩手中、现在正在自由下落的是一个玩偶。

“去吧，爱丽丝。”玲听到女孩小声地说着。

“别想不开啊，”玲站在女孩身后，气喘吁吁地开口说，“你还这么年轻。我们不是还有这个美丽的世界、美好的每一天吗？”这话说出口，玲立刻后悔了；这样的话是万万不能对一个即将寻短见的人说的。玲不知道怎么办，只能尴尬地站在原地。

不对，如果第二种假说成立的话，玩偶应该向上下落而不是向下下落。还是不对，如果地球的引力场已经受到这种程度的扭曲的话，地球早就已经超过洛希极限而解体了……停停，现在是想这个的时候吗？

女孩转过身来。她轻盈地越过栏杆，就像猫一样，面对手足无措的玲。玲发现女孩的脸上是灿烂的笑容。

“你来啦，爱丽丝。”

“爱丽丝？谁？谁是爱丽丝？你是说你扔下去的玩偶吗？”

女孩轻轻地摇了摇头，“走到这一步，来到这里，正是爱丽丝你的意愿吧。”

她有回答我的问题吗？玲不解，便不再纠结这个问题了，“你叫什么名字？我怎么称呼你？”

“雪莉，”女孩坚定地说，“我叫雪莉。你不记得我了吗？”

虽然玲略微有点脸盲，尤其是对孩子，但她对雪莉完全没有印象，“我之前见过你吗？”

雪莉没有回答。看来谈话又进入了死胡同；玲只好换个话题：“你刚才在做什么？”

“我在研究空气力学哦。”

空气力学？没听过。大概是什么新兴的交叉学科吧。听起来像是通过研究空气的运动来预测天气和气候的变化。“这肯定需要个参数量不小的模型来做预测。”玲心想。

“你真的没有打算跳下去吧？”玲不确定地问道。

“怎么会呢，”雪莉笑了笑，但她的笑脸上流下了眼泪，“你终于来了。太好了。”接着，她拿出手机，“我要把这个好消息告诉所有人。今天是雪莉和爱丽丝的幸运日哦。”

“所有人？”

“没错，”雪莉按着手机，“通讯录里的所有人。群发哦！”

玲摸不着头脑。她不知道雪莉为何认识自己，也不理解雪莉究竟在干什么。“如果用精神分裂症来解释的话，那就一切都说得通了。”玲想着，“不过，这种桥段即使放在小说里，也会被读者喷的。”

“你接下来要去天空的另一边吗？”看样子雪莉已经发完了讯息，收起了手机；但她旋即问出更令人不解的一句。

“什么意思？”

雪莉狡黠地一笑，“你会理解的，最终。”

玲决定顺着她的话，“那你和我一起去吗？”

“不哦，爱丽丝只能自己去。”雪莉注视着玲，“但是雪莉会祝福你的。愿雪莉的祝福保佑你。”

雪莉说完便消失不见了。玲抬头，只见猩红的天空向她扑来。

\includepdf[pages=-]{supp1.pdf}
